\section{Introduction}
\label{sec:intro}

Broken access control heads the OWASP Top~10~\cite{owasp2021a01}. Most
research tools for it ask whether the \emph{server} enforces authorization,
by comparing or replaying requests under different identities
\cite{sun2011static,li2014blackbox,noseevich2011detecting,monshizadeh2014mace,liu2025bacscan}.
A related class of defect lives one layer up. Consider a dashboard that
renders an \emph{Admin} link for every role. The backend correctly rejects a
guest's request, so a request-level scanner reports nothing. The link itself
still leaks the existence and location of privileged functionality, invites
probing, and violates the organisation's least-privilege UI policy. The same
applies to a debug panel left in a production build, a destructive control a
guest can reach in two clicks, or the \emph{absence} of an element the policy
requires, such as a skip-to-content link, a consent banner, or a logout
affordance.

Workflow-based UI testing does not catch these defects, because tests check
what \emph{should} happen and nobody writes a test for a button that should
not exist. Model-inference crawlers~\cite{mesbah2012crawljax,memon2003ripping,dallmeier2013webmate}
explore autonomously but check generic or hand-coded invariants. DOM
invariants~\cite{mesbah2012invariant} can express presence and absence, but as
per-page XPath predicates written in code, with no notion of role. LLM and
VLM GUI agents~\cite{liu2024gptdroid,yoon2024droidagent,liu2024visiondroid,zheng2024seeact}
explore or judge screens in natural language, but they let a model issue the
verdict, which is risky when the verdict is ``this element is \emph{not}
there'': VLMs are known to handle negation poorly~\cite{alhamoud2025negation}.

\paragraph{This paper.}
We present \sys{}, which reframes UI testing as \emph{state-space auditing}.
Given a seed URL, a role, and a policy of the form ``a guest must never see the
admin link; every page must offer a skip-to-content link'', \sys{} explores the
application autonomously, keeps the full cyclic state graph, and audits every
reached state. Each finding carries the policy clause it breaks, the set
operation that caught it, and human-checkable evidence.

\paragraph{Contributions.}
We do not claim to be the first to crawl web applications into state graphs,
to check invariants on crawled states, or to use VLMs for GUI testing; each has
substantial prior art (\cref{sec:related}). Our contributions are a
formalisation and a design that \emph{combine} these threads for a question
none of them targets: \emph{does the UI a given role can reach conform to a
declared policy of what must and must not exist?}
\begin{enumerate}
  \item \textbf{A polarity-typed, role-scoped policy oracle}
    (\cref{sec:model}). Forbidden-present ($S\cap\mathcal F$) and
    required-absent ($\mathcal R\setminus S$) are distinct set operations over
    a \emph{semantic} state abstraction (roles, accessible names, tags,
    capabilities), not raw XPath. Polarity is part of the policy schema, and
    every violation is tagged with its clause type.
  \item \textbf{An error-direction asymmetry result}
    (\cref{prop:asym}) together with its empirical confirmation
    (\cref{sec:eval:rq2}). Perception errors propagate in opposite directions
    for the two clause types. This argues for reporting accuracy \emph{per
    polarity} and for a \emph{negation-safe division of labour}, in which
    deterministic DOM/AX evidence decides existence and a VLM only annotates.
  \item \textbf{A side-effect-aware cyclic exploration loop}
    (\cref{sec:system}). It uses claim-before-execute per $(v,a)$ pair,
    kind-restricted replay, content-addressed actions, and normalized
    fingerprints, and runs as a checkpointable LangGraph state machine behind
    typed ports.
  \item \textbf{An honest, reproducible evaluation} on a polarity-complete
    ground-truth application (\cref{sec:eval}). It separates what the loop
    guarantees (exact recovery with an oracle perceiver) from what perception
    currently delivers, and localizes the remaining gap to semantic subject
    grounding.
\end{enumerate}

Unlike access-control scanners, \sys{} does not test server-side enforcement.
It audits what each role \emph{sees}, so the two are complementary.
